% =====================================================================
%  Bölüm 6 — Yüklemler Mantığı: ∀, ∃ ve Eşitlik
% =====================================================================
\chapter{Yüklemler Mantığı: \texorpdfstring{$\forall$}{∀}, \texorpdfstring{$\exists$}{∃} ve Eşitlik}\label{bol:yuklemler}

\begin{bolumhedefleri}
\begin{itemize}
  \item Yüklemleri \emph{önerme değerli fonksiyonlar} olarak görmeyi ve
        \Lean{}'de \lc{α → Prop} türüyle yazmayı,
  \item $\forall$ ve $\exists$ için giriş ve eleme kurallarını, özellikle
        \emph{öz değişken koşulunu},
  \item Niceleyicilerin değilleme ve sırayla nasıl etkileştiğini; hangi
        etkileşimlerin yapıcı, hangilerinin klasik olduğunu,
  \item Serbest ve bağlı değişkenleri, yerine koymayı ve değişken
        yakalama tuzağını,
  \item Eşitliği: \lc{rfl}, \leantac{rw}, \lc{▸}, \lc{calc} ve Leibniz ilkesini,
  \item Berber paradoksunu \Lean{}'de ispatlamayı ve boş evren inceliğini,
  \item Curry--Howard tablosunun bağımlı tiplere genişlemesini:
        $\forall = \Pi$, $\exists = \Sigma$.
\end{itemize}
\end{bolumhedefleri}

Önermeler mantığı, ``Sokrates ölümlüdür'' cümlesini bölünmez bir $P$
harfi olarak görür. Oysa Bölüm~1'deki kıyasın geçerliliği tam da bu
cümlenin \emph{içindeki} yapıya dayanıyordu: ``Bütün insanlar
ölümlüdür; Sokrates bir insandır; o hâlde Sokrates ölümlüdür.'' Bu
çıkarımı ifade etmek için önermelerin içine bakmamız, \emph{bireyler} ve
onların \emph{özellikleri} hakkında konuşmamız gerekir. Bu bölümün konusu
budur.

% ---------------------------------------------------------------------
\section{Önermelerden yüklemlere}\label{sec:yuklemler}

\begin{tanim}{Yüklem}{yuklem}
Bir $\alpha$ \textbf{tanım evreni} üzerindeki bir \textbf{yüklem},
$\alpha$'nın her elemanına bir önerme karşılık getiren bir
fonksiyondur:
\[
  P : \alpha \to \mathsf{Prop}.
\]
$P(a)$, ``$a$, $P$ özelliğini taşır'' diye okunur. İki argümanlı bir
yüklem ($R : \alpha \to \beta \to \mathsf{Prop}$) bir \textbf{ilişkidir}.
\end{tanim}

\Lean{}'de yüklemler sıradan fonksiyonlardır. ``Çifttir'' yüklemi, bir
doğal sayı alıp bir önerme döndüren fonksiyondur:

\leandosya[Yüklemlerin türleri]{Bolum06}{yuklem_tip}

\begin{figure}[htbp]
\centering
\begin{tikzpicture}[
    eleman/.style={circle, fill=anarenk, inner sep=1.8pt},
    onerme/.style={draw=#1, fill=#1!8, rounded corners=3pt, font=\small,
                   inner sep=4pt, minimum width=2.2cm},
    ok/.style={-{Stealth[length=2.2mm]}, thick, black!45},
  ]
  \draw[anarenk, thick, fill=tanimbg, rounded corners=10pt] (-0.9,-2.6) rectangle (0.9,1.0);
  \node[font=\baslikfont\bfseries\small, text=anarenk] at (0,1.35) {$\alpha = \mathbb{N}$};
  \foreach \n/\y in {0/0.5, 1/-0.2, 2/-0.9, 3/-1.6} {
    \node[font=\small] (n\n) at (0,\y) {$\n$};
  }
  \node[font=\small] at (0,-2.2) {$\vdots$};
  \node[onerme=bittiframe] (p0) at (4.6,0.5)  {$0 \% 2 = 0$};
  \node[onerme=hataframe]  (p1) at (4.6,-0.2) {$1 \% 2 = 0$};
  \node[onerme=bittiframe] (p2) at (4.6,-0.9) {$2 \% 2 = 0$};
  \node[onerme=hataframe]  (p3) at (4.6,-1.6) {$3 \% 2 = 0$};
  \foreach \n in {0,1,2,3} \draw[ok] (n\n) -- (p\n);
  \draw[ok, dashed] (0.3,-2.2) -- (3.4,-2.2);
  \node[font=\baslikfont\bfseries\small, text=vurgu] at (2.1,1.35) {$P$};
  \node[font=\baslikfont\bfseries\small, text=koyu] at (4.6,1.35) {$\mathsf{Prop}$};
  \node[font=\baslikfont\scriptsize, text=bittiframe, anchor=west] at (6.0,0.15) {doğru (kanıtı var)};
  \node[font=\baslikfont\scriptsize, text=hataframe, anchor=west] at (6.0,-0.55) {yanlış (kanıtı yok)};
\end{tikzpicture}
\caption{``Çifttir'' yüklemi $P(n) := (n \% 2 = 0)$, her doğal sayıya
bir önerme karşılık getirir. Önermelerin bazıları doğru, bazıları
yanlıştır; ama hepsi birer önermedir.}
\label{fig:yuklem}
\end{figure}

Artık Bölüm~1'deki kıyası tam olarak yazabiliriz:

\leandosya[Sokrates kıyası]{Bolum06}{sokrates}

\noindent İspatın tamamı \lc{h1 sokrates h2}'dir. \lc{h1}'i önce
\lc{sokrates}'e, sonra \lc{h2}'ye uyguluyoruz. ``Bütün insanlar
ölümlüdür'' önermesinin kanıtı, \emph{her birey için} ``insansa
ölümlüdür'' kanıtını üreten bir fonksiyondur.

% ---------------------------------------------------------------------
\section{Evrensel niceleyici (\texorpdfstring{$\forall$}{∀})}\label{sec:forall}

\begin{center}
\begin{tabular}{c@{\hspace{5em}}c}
\begin{prooftree}
  \hypo{\mathcal{D}}
  \ellipsis{}{P(a)}
  \infer1[$\forall$G$^{*}$]{\forall x,\; P(x)}
\end{prooftree}
&
\begin{prooftree}
  \hypo{\forall x,\; P(x)}
  \infer1[$\forall$E]{P(t)}
\end{prooftree}
\\[10pt]
\small\textcolor{ikincil}{\bfseries giriş:} \leantac{intro x} &
\small\textcolor{vurgu}{\bfseries eleme:} \lc{h t}, \leantac{specialize h t}
\end{tabular}
\end{center}

\noindent {\small $^{*}$\,\textbf{Öz değişken koşulu:} $a$, $\mathcal{D}$
türetiminin açık varsayımlarında geçmemelidir.}

\textbf{$\forall$ girişi.} ``Her $x$ için $P(x)$'' göstermek için
\emph{keyfi} bir $a$ seçer ve $P(a)$'yı gösteririz. Matematikçilerin
``keyfi bir $x$ alalım'' cümlesi budur. Kilit sözcük \emph{keyfi}dir:
$a$ hakkında hiçbir özel varsayım yapmamış olmalıyız. Aksi hâlde
``$a = 2$ olsun; $a$ çifttir; o hâlde her sayı çifttir'' gibi bir saçmalık
türetirdik. Bu kısıta \textbf{öz değişken koşulu} (İng.\
\emph{eigenvariable condition}) denir.

\textbf{$\forall$ elemesi.} ``Her $x$ için $P(x)$'' biliyorsak, herhangi bir
$t$ terimi için $P(t)$'yi biliriz. Bu, ``genelden özele'' geçiştir.

\leandosya[$\forall$ girişi]{Bolum06}{forall_giris}

\begin{hedefgecisi}
\begin{hedefdurumu}[Önce]
α : Type
⊢ ∀ (x : α), x = x
\end{hedefdurumu}
\sonra
\begin{hedefdurumu}[\texttt{intro x} sonrası]
α : Type
x : α
⊢ x = x
\end{hedefdurumu}
\end{hedefgecisi}

\noindent \leantac{intro x}, bağlama yeni bir \lc{x : α} ekler. \Lean{},
öz değişken koşulunu kendiliğinden sağlar: \lc{x} yeni bir addır ve
bağlamdaki hiçbir hipotez onun hakkında bir şey söylemez.

\begin{onemlifikir}[Önemli Fikir: $\forall$ bir fonksiyondur, $\to$ onun özel hâlidir]
\leantac{intro}'nun hem $\to$ hem $\forall$ için çalışması bir rastlantı
değildir. \Lean{}'de \lc{∀ x : α, P x} türünde bir kanıt, her \lc{a : α}
için \lc{P a} türünde bir kanıt üreten bir fonksiyondur. Tek fark,
\emph{dönüş türünün argümana bağlı olmasıdır}: \lc{h 3 : P 3},
\lc{h 5 : P 5}. Böyle fonksiyonlara \textbf{bağımlı fonksiyon} denir.
\lc{P → Q} ise dönüş türü argümana bağlı olmayan özel bir durumdur.
\end{onemlifikir}

\leandosya[$\forall$ elemesi]{Bolum06}{forall_eleme}

\noindent İki kuralı birlikte kullanan bir örnek:

\leandosya[$\forall$ ve $\wedge$]{Bolum06}{forall_ve}

\noindent Bazen bir genel hipotezi belirli bir bireye ``indirgemek''
isteriz. \leantac{specialize} taktiği, hipotezi o bireye uygulanmış
hâliyle değiştirir:

\leandosya[\texttt{specialize}]{Bolum06}{specialize}

% ---------------------------------------------------------------------
\section{Varoluşsal niceleyici (\texorpdfstring{$\exists$}{∃})}\label{sec:exists}

\begin{center}
\begin{tabular}{c@{\hspace{5em}}c}
\begin{prooftree}
  \hypo{P(t)}
  \infer1[$\exists$G]{\exists x,\; P(x)}
\end{prooftree}
&
\begin{prooftree}
  \hypo{\exists x,\; P(x)}
  \hypo{[P(a)]^{n}}
  \ellipsis{}{R}
  \infer2[$\exists$E$^{n,*}$]{R}
\end{prooftree}
\\[10pt]
\small\textcolor{ikincil}{\bfseries giriş:} \lc{⟨t, h⟩} &
\small\textcolor{vurgu}{\bfseries eleme:} \leantac{obtain ⟨a, ha⟩ := h}
\end{tabular}
\end{center}

\noindent {\small $^{*}$\,\textbf{Öz değişken koşulu:} $a$, $R$'de ve
$P(a)$ dışındaki açık varsayımlarda geçmemelidir.}

\textbf{$\exists$ girişi.} ``$P$'yi sağlayan bir $x$ vardır'' göstermek
için somut bir \textbf{tanık} $t$ verir ve $P(t)$'yi gösteririz.
BHK yorumuyla uyumlu olarak, bir varlık kanıtı \emph{tanığını taşır}.

\textbf{$\exists$ elemesi.} ``$P$'yi sağlayan bir $x$ vardır'' biliyorsak,
ona bir ad veririz: ``Öyle bir $a$ olsun ki $P(a)$.'' Bu $a$ hakkında
$P(a)$'dan başka hiçbir şey bilmeyiz. Öz değişken koşulu, $a$'nın
sonuca ``sızmasını'' engeller: $a$ yalnızca ispatın içinde yaşayan geçici
bir addır.

\leandosya[$\exists$ girişi]{Bolum06}{exists_giris}

\noindent Köşeli ayraçlar, \lc{∧}'deki gibi bir çift kurar: Birinci
bileşen tanık, ikinci bileşen tanığın özelliğinin kanıtıdır. İkinci
ispatta \leantac{refine} önce tanığı verir; geriye kalan hedef \lc{4 > 3}
olur.

\leandosya[$\exists$ elemesi]{Bolum06}{exists_eleme}

\begin{hedefgecisi}
\begin{hedefdurumu}[\texttt{obtain} öncesi]
h1 : ∀ (x : α), P x → Q x
h2 : ∃ x, P x
⊢ ∃ x, Q x
\end{hedefdurumu}
\sonra
\begin{hedefdurumu}[\texttt{obtain} sonrası]
h1 : ∀ (x : α), P x → Q x
a : α
ha : P a
⊢ ∃ x, Q x
\end{hedefdurumu}
\end{hedefgecisi}

\noindent Terim modunda aynı ispat, \leankw{let} ile örüntü eşleme
kullanır:

\leandosya[$\exists$ elemesi, terim modu]{Bolum06}{exists_eleme_terim}

\begin{figure}[htbp]
\centering
\begin{tikzpicture}[
    evren/.style={draw=#1, thick, fill=#1!6, rounded corners=12pt,
                  minimum width=4.2cm, minimum height=2.8cm},
    nokta/.style={circle, inner sep=2.4pt, fill=#1},
  ]
  % ∀
  \node[evren=ikincil] (A) at (0,0) {};
  \foreach \x/\y in {-1.3/0.7, -0.4/0.9, 0.6/0.6, 1.4/0.8, -1.1/-0.3, 0.0/-0.1,
                     1.0/-0.4, -0.5/-0.9, 0.5/-1.0, 1.5/-0.9}
    \node[nokta=bittiframe] at (\x,\y) {};
  \node[font=\baslikfont\bfseries, text=ikincil] at (0,1.75) {$\forall x,\; P(x)$};
  \node[font=\baslikfont\scriptsize, text=black!60, align=center] at (0,-1.85)
       {\emph{her} eleman $P$'yi sağlar\\kanıt: her eleman için bir kanıt (fonksiyon)};
  % ∃
  \node[evren=vurgu] (B) at (6.4,0) {};
  \foreach \x/\y in {-1.3/0.7, -0.4/0.9, 0.6/0.6, 1.4/0.8, -1.1/-0.3,
                     1.0/-0.4, -0.5/-0.9, 0.5/-1.0, 1.5/-0.9}
    \node[nokta=black!25] at ({6.4+\x},\y) {};
  \node[nokta=bittiframe, minimum size=9pt] (t) at (6.4,-0.1) {};
  \node[font=\baslikfont\scriptsize\bfseries, text=bittiframe, right=1pt of t] {tanık $t$};
  \node[font=\baslikfont\bfseries, text=vurgu] at (6.4,1.75) {$\exists x,\; P(x)$};
  \node[font=\baslikfont\scriptsize, text=black!60, align=center] at (6.4,-1.85)
       {\emph{en az bir} eleman $P$'yi sağlar\\kanıt: bir tanık ve onun kanıtı (çift)};
\end{tikzpicture}
\caption{İki niceleyicinin anlamı ve BHK yorumuna göre kanıtları.}
\label{fig:niceleyiciler}
\end{figure}

% ---------------------------------------------------------------------
\section{Niceleyicilerin etkileşimi}\label{sec:niceleyici-etkilesim}

\subsection{Değilleme ve niceleyiciler}

Niceleyiciler değillemeyle De Morgan yasalarına benzer biçimde etkileşir.
Hatta $\forall$, büyük bir ``ve'', $\exists$ ise büyük bir ``veya''
gibidir: Evren $\{a, b\}$ ise $\forall x\,P(x) = P(a) \wedge P(b)$ ve
$\exists x\,P(x) = P(a) \vee P(b)$ olur. Bölüm~3'teki asimetri burada
da karşımıza çıkar:

\leandosya[``Hiç kimse'' = ``herkes değil'']{Bolum06}{degil_exists}
\leandosya[Yapıcı yön]{Bolum06}{degil_forall_kolay}
\leandosya[Klasik yön]{Bolum06}{degil_forall_klasik}

\noindent Son ispatın neden klasik mantık gerektirdiğini BHK yorumuyla
anlayabiliriz. Hedef bir $\exists$ önermesi olduğuna göre bir tanık
vermemiz gerekir. Elimizdeki tek bilgi ise ``herkesin $P$'yi sağladığı
doğru değildir''dir. Bu bilgi, $P$'yi sağlamayan bireyin \emph{kim}
olduğunu söylemez.

\begin{table}[htbp]
\centering
\renewcommand{\arraystretch}{1.5}
\rowcolors{2}{tablosatir}{white}
\begin{tabular}{>{\centering\arraybackslash}p{5.6cm}
                >{\centering\arraybackslash}p{2.0cm}
                >{\centering\arraybackslash}p{2.0cm}
                >{\raggedright\arraybackslash}p{2.6cm}}
\rowcolor{tabloust}
\tablobaslik{Önerme} & \tablobaslik{$\to$} & \tablobaslik{$\leftarrow$} & \tablobaslik{Lean} \\
$\neg\exists x\,P(x) \leftrightarrow \forall x\,\neg P(x)$ & \textcolor{bittiframe}{yapıcı} & \textcolor{bittiframe}{yapıcı} & \lc{degil\_exists} \\
$\neg\forall x\,P(x) \leftrightarrow \exists x\,\neg P(x)$ & \textcolor{hataframe}{klasik} & \textcolor{bittiframe}{yapıcı} & \lc{degil\_forall} \\
$\forall x\,(P \wedge Q) \leftrightarrow \forall x\,P \wedge \forall x\,Q$ & \textcolor{bittiframe}{yapıcı} & \textcolor{bittiframe}{yapıcı} & Alıştırma~6.4 \\
$\exists x\,(P \vee Q) \leftrightarrow \exists x\,P \vee \exists x\,Q$ & \textcolor{bittiframe}{yapıcı} & \textcolor{bittiframe}{yapıcı} & Alıştırma~6.5 \\
$\exists x\,(P \wedge Q) \to \exists x\,P \wedge \exists x\,Q$ & \textcolor{bittiframe}{yapıcı} & \textcolor{hataframe}{yanlış} & Alıştırma~6.6 \\
$\exists x\,\forall y\,R \to \forall y\,\exists x\,R$ & \textcolor{bittiframe}{yapıcı} & \textcolor{hataframe}{yanlış} & \lc{exists\_forall\_sira} \\
\end{tabular}
\caption{Niceleyicilerin etkileşimleri. ``Yanlış'', o yönün genel olarak
geçerli olmadığı, yani bir karşı örneği bulunduğu anlamına gelir.}
\label{tab:niceleyici-etkilesim}
\end{table}

\subsection{Niceleyicilerin sırası}

Niceleyicilerin sırası anlamı kökten değiştirir:
\begin{center}
\begin{tabular}{rl}
$\forall y\,\exists x\; \mathit{Sever}(x, y)$: & ``Herkesi seven \emph{biri} vardır'' \emph{değil}, ``herkesin bir seveni vardır.'' \\
$\exists x\,\forall y\; \mathit{Sever}(x, y)$: & ``Herkesi seven biri vardır.''
\end{tabular}
\end{center}

\noindent İkincisi birincisinden güçlüdür:

\leandosya[$\exists\forall \Rightarrow \forall\exists$]{Bolum06}{sira}

\noindent Ters yön yanlıştır. Doğal sayılarda ``her sayıdan büyük bir sayı
vardır'' doğrudur, ama ``her sayıdan büyük bir sayı'' yoktur
(\cref{fig:sira}):

\leandosya[Karşı örnek]{Bolum06}{sira_karsi}

\begin{figure}[htbp]
\centering
\begin{tikzpicture}[
    sayi/.style={circle, draw=anarenk, fill=tanimbg, thick, minimum size=0.75cm,
                 font=\small},
    ok/.style={-{Stealth[length=2.2mm]}, thick, #1},
  ]
  \node[font=\baslikfont\bfseries\small, text=bittiframe] at (2.5,1.3) {$\forall y\,\exists x,\; x > y$ \quad (doğru)};
  \foreach \i in {0,...,4} {
    \node[sayi] (y\i) at ({1.25*\i},0) {$\i$};
    \node[sayi, draw=bittiframe, fill=bittibg] (x\i) at ({1.25*\i},-1.6) {$\pgfmathparse{int(\i+1)}\pgfmathresult$};
    \draw[ok=bittiframe] (y\i) -- (x\i);
  }
  \node[font=\small] at (6.25,0) {$\cdots$};
  \node[font=\baslikfont\scriptsize, text=black!60] at (2.5,-2.4) {her $y$ için ayrı bir tanık: $x = y+1$};

  \node[font=\baslikfont\bfseries\small, text=hataframe] at (10.3,1.3) {$\exists x\,\forall y,\; x > y$ \quad (yanlış)};
  \node[sayi, draw=hataframe, fill=hatabg] (X) at (10.3,-1.6) {$x$};
  \foreach \i in {0,...,3} {
    \node[sayi] (z\i) at ({8.4+1.25*\i},0) {};
    \draw[ok=hataframe] (X) -- (z\i);
  }
  \node[font=\baslikfont\scriptsize, text=hataframe, align=center] at (10.3,-2.5)
       {tek bir $x$ hepsinden büyük olamaz:\\$y = x$ seçilince $x > x$ gerekir};
\end{tikzpicture}
\caption{Niceleyicilerin sırası. $\forall\exists$'de tanık $y$'ye bağlı
olabilir; $\exists\forall$'de tek bir tanık herkese yetmelidir.}
\label{fig:sira}
\end{figure}

\begin{dikkat}[Dikkat: Tanık neye bağlıdır?]
$\forall y\,\exists x$ önermesinin kanıtı, Curry--Howard altında, her
$y$ için bir $x$ \emph{üreten bir fonksiyondur}: \lc{fun y => ⟨y + 1, \ldots⟩}.
Tanık $y$'ye bağlıdır. $\exists x\,\forall y$ ise tek bir sabit tanık
ister. Analizdeki ``sürekli'' ile ``düzgün sürekli'' arasındaki farkın
özü de budur: $\delta$, $x$ noktasına bağlı olabilir mi?
\end{dikkat}

% ---------------------------------------------------------------------
\section{Serbest ve bağlı değişkenler}\label{sec:serbest-bagli}

$\forall x,\; x + y = y + x$ önermesinde $x$ ve $y$ değişkenlerinin
rolleri farklıdır. $x$ niceleyiciye \emph{bağlıdır}: Onun yerine başka
bir ad koysak ($\forall z,\; z + y = y + z$) önermenin anlamı değişmez.
$y$ ise \emph{serbesttir}: Önermenin doğruluğu $y$'nin ne olduğuna
bağlıdır.

\begin{tanim}{Serbest ve bağlı değişken}{serbest-bagli}
Bir önermede bir $x$ değişkeninin bir geçişi, bir $\forall x$ ya da
$\exists x$ niceleyicisinin kapsamı içindeyse \textbf{bağlı}, değilse
\textbf{serbesttir}. Bağlı değişkenlerin adları değiştirilerek elde
edilen önermelere \textbf{$\alpha$-denk} önermeler denir ve aynı önerme
sayılırlar.
\end{tanim}

\paragraph{Yerine koyma ve değişken yakalama}
$\forall$E kuralı $\forall x\,P(x)$'ten $P(t)$'ye geçer: $x$'in yerine $t$
terimini koyar. Burada ince bir tuzak vardır (\cref{fig:yakalama}).

\begin{figure}[htbp]
\centering
\begin{tikzpicture}[
    kutu/.style={draw=#1, fill=#1!7, thick, rounded corners=4pt,
                 text width=5.6cm, align=center, inner sep=6pt, font=\small},
    ok/.style={-{Stealth[length=2.6mm]}, very thick, #1},
  ]
  \node[kutu=anarenk] (a) at (0,0)
    {$\forall x,\; \exists y,\; y > x$\\[2pt]{\scriptsize ``Her sayıdan büyük bir sayı vardır.'' (doğru)}};
  \node[kutu=hataframe] (b) at (-3.4,-2.6)
    {$\exists y,\; y > \textcolor{hataframe}{\mathbf{y}}$\\[2pt]{\scriptsize $x$ yerine $y$ konuldu; $y$ \emph{yakalandı}. (yanlış!)}};
  \node[kutu=bittiframe] (c) at (3.4,-2.6)
    {$\exists z,\; z > \textcolor{bittiframe}{\mathbf{y}}$\\[2pt]{\scriptsize Önce bağlı $y$, $z$ yapıldı; sonra yerine koyma. (doğru)}};
  \draw[ok=hataframe] (a) -- node[left=6pt, font=\baslikfont\scriptsize, text=hataframe]{saf yerine koyma} (b);
  \draw[ok=bittiframe] (a) -- node[right=6pt, font=\baslikfont\scriptsize, text=bittiframe]{yakalamadan kaçınan} (c);
\end{tikzpicture}
\caption{Değişken yakalama: $\forall$E ile $x$ yerine $y$ terimini koyarken
içerideki $\exists y$ niceleyicisi, dışarıdan gelen $y$'yi ``yakalar''.}
\label{fig:yakalama}
\end{figure}

Doğru yerine koyma, gerektiğinde bağlı değişkenlerin adlarını değiştirir.
Kâğıt üzerinde bu dikkat gerektiren bir iştir ve mantık kitaplarının
tanımlarında sık rastlanan hataların kaynağıdır. \Lean{} ise bağlı
değişkenleri içeride ad yerine \emph{de Bruijn indeksleriyle} (``kaç
niceleyici yukarıdaki değişken'') tutar; böylece yakalama hiç oluşmaz.
Bu temsili \cref{bol:sozdizim}'te kendi formül tipimizi tanımlarken
kullanacağız.

% ---------------------------------------------------------------------
\section{Eşitlik}\label{sec:esitlik}

Eşitlik, yüklemler mantığının özel bir ilişkisidir. Kuralları şöyledir:

\begin{center}
\begin{tabular}{c@{\hspace{5em}}c}
\begin{prooftree}
  \infer0[$=$G]{t = t}
\end{prooftree}
&
\begin{prooftree}
  \hypo{s = t}
  \hypo{P(s)}
  \infer2[$=$E]{P(t)}
\end{prooftree}
\\[10pt]
\small\textcolor{ikincil}{\bfseries giriş (yansıma):} \lc{rfl} &
\small\textcolor{vurgu}{\bfseries eleme (yerine koyma):} \leantac{rw}, \lc{h ▸ hp}
\end{tabular}
\end{center}

\noindent Yalnızca bu iki kuraldan eşitliğin bütün bilinen özellikleri
türer: simetri, geçişlilik, fonksiyonlarla uyum.

\leandosya[Simetri ve geçişlilik]{Bolum06}{esitlik_temel}
\leandosya[Yeniden yazma]{Bolum06}{esitlik_rw}

\noindent \leantac{rw [h]}, \lc{h : a = b} eşitliğini kullanarak hedefteki
her \lc{a}'yı \lc{b} yapar. Yeniden yazdıktan sonra hedef \lc{c + 1 = c + 1}
gibi apaçık bir eşitliğe dönüşürse, \leantac{rw} onu \lc{rfl} ile
kendiliğinden kapatır. Ters yönde yazmak için \lc{rw [← h]} kullanılır.

\leandosya[Yerine koyma]{Bolum06}{esitlik_subst}
\leandosya[Fonksiyonlarla uyum]{Bolum06}{esitlik_congr}

\paragraph{calc: eşitlik zincirleri}
Matematikte uzun hesapları bir eşitlik zinciri olarak yazarız:
$a = b = c = d$. \leankw{calc} bloğu bunun \Lean{}'deki karşılığıdır.
Her satır bir adımı ve gerekçesini verir; alt çizgi bir önceki satırın
sağ tarafını gösterir:

\leandosya[Eşitlik zinciri]{Bolum06}{calc_ornek}

\noindent \leankw{calc}, \cref{bol:tumevarim}'de aritmetik ispatlarda ve
\lc{<}, \lc{≤} gibi ilişkilerin zincirlerinde çok işimize yarayacak.

\paragraph{Leibniz ilkesi}
Gottfried Leibniz'e göre iki şey, bütün özellikleri ortaksa aynıdır.
Bu ilke, eşitliğin $=$E kuralının tersidir ve \Lean{}'de tek satırla
ispatlanır:

\leandosya[Leibniz ilkesi]{Bolum06}{leibniz}

\noindent İspattaki hile, ``$a$'ya eşit olmak'' özelliğini seçmektir:
$a$ bu özelliği \lc{rfl} ile taşır, o hâlde $b$ de taşır, yani $a = b$.

% ---------------------------------------------------------------------
\section{Berber paradoksu}\label{sec:berber}

\begin{gunlukornek}[Bir köyün berberi]
Bir köyde bir berber, ``kendini tıraş etmeyen herkesi tıraş eder,
kendini tıraş edeni ise tıraş etmez.'' Berber kendini tıraş eder mi?
\begin{itemize}
  \item Ederse, kendini tıraş eden biri olur; o hâlde onu tıraş etmez.
  \item Etmezse, kendini tıraş etmeyen biri olur; o hâlde onu tıraş eder.
\end{itemize}
Her iki durumda da çelişki çıkar. Sonuç: Böyle bir berber \emph{yoktur}.
\end{gunlukornek}

\begin{figure}[htbp]
\centering
\begin{tikzpicture}[
    kisi/.style={circle, draw=anarenk, fill=tanimbg, thick, minimum size=0.9cm,
                 font=\small},
    ok/.style={-{Stealth[length=2.4mm]}, thick, #1},
  ]
  \node[kisi, draw=vurgu, fill=teorembg, minimum size=1.1cm] (b) at (0,0) {$b$};
  \node[kisi] (a) at (-2.6,1.2) {$a$};
  \node[kisi] (c) at (2.6,1.2) {$c$};
  \node[kisi] (d) at (-2.6,-1.2) {$d$};
  \node[kisi] (e) at (2.6,-1.2) {$e$};
  % kendini tıraş edenler: a, c (döngü)
  \draw[ok=black!45] (a) to[loop above, looseness=6] (a);
  \draw[ok=black!45] (c) to[loop above, looseness=6] (c);
  % berber d ve e'yi tıraş eder
  \draw[ok=vurgu] (b) -- (d);
  \draw[ok=vurgu] (b) -- (e);
  \draw[ok=vurgu, dashed] (b) to[loop below, looseness=6] node[below, font=\baslikfont\large\bfseries, text=vurgu]{?} (b);
  \node[font=\baslikfont\scriptsize, text=black!60, align=left, anchor=west] at (4.0,0.6)
       {$a$, $c$: kendini tıraş eder\\$\Rightarrow$ berber onları tıraş etmez};
  \node[font=\baslikfont\scriptsize, text=vurgu, align=left, anchor=west] at (4.0,-0.6)
       {$d$, $e$: kendini tıraş etmez\\$\Rightarrow$ berber onları tıraş eder};
\end{tikzpicture}
\caption{Berber $b$'nin kendine giden oku çizilmeli midir? Çizilse de
çizilmese de kural bozulur.}
\label{fig:berber}
\end{figure}

\leandosya[Berber paradoksu]{Bolum06}{berber}

\noindent İspatın çekirdeği Alıştırma~3.8 ile aynıdır: Hiçbir önerme
kendi değillemesine denk olamaz. \lc{hb b} bize
\lc{TirasEder b b ↔ ¬ TirasEder b b} denkliğini verir; geri kalanı
önermeler mantığıdır. Dikkat edin, ispat \emph{yapıcıdır}: Paradoks,
klasik mantığa hiç ihtiyaç duymaz.

\begin{tarihce}[Tarihçe: Russell'ın mektubu]
Berber hikâyesi, Bertrand Russell'ın 1901'de bulduğu ve 1902'de Frege'ye
bir mektupla bildirdiği paradoksun popüler bir anlatımıdır. Russell'ın
asıl sorusu şuydu: ``Kendini eleman olarak içermeyen bütün kümelerin
kümesi'' kendini içerir mi? Frege'nin aritmetiği mantıktan türetme
programı, sınırsız küme oluşturmaya izin veren bir ilkeye dayandığı için
bu paradoksla çöktü \citep{vanheijenoort1967}. Paradoksa verilen
cevaplardan biri Russell'ın \emph{tip kuramıydı}: Bir kümenin elemanları
ile kümenin kendisi farklı ``tiplerde'' olmalıdır. \Lean{}'in tip
sistemi, bu fikrin yüz yıllık evriminin bir ürünüdür.
\end{tarihce}

% ---------------------------------------------------------------------
\section{Boş evren inceliği}\label{sec:bos-evren}

Klasik mantık kitaplarının çoğunda şu kural geçerlidir:
$\forall x\,P(x) \to \exists x\,P(x)$. ``Herkes $P$ ise, $P$ olan biri
vardır.'' Kulağa apaçık geliyor, ama evrende \emph{hiç kimse yoksa}
ne olur?

\leandosya[Boş ve boş olmayan evrenler]{Bolum06}{bos_evren}

\noindent Boş tip \lc{Empty} üzerinde $\forall x\,P(x)$ \emph{boş yere
doğrudur}: Kontrol edilecek hiç eleman yoktur. Ama $\exists x\,P(x)$
yanlıştır: Tanık bulunamaz. Dolayısıyla $\forall \to \exists$ geçişi,
yalnızca evren boş değilse geçerlidir. \Lean{}'de bu varsayım
\lc{[Inhabited α]} ile açıkça yazılır.

\begin{dikkat}[Dikkat: Geleneksel mantık evreni boş olmayan kabul eder]
Birinci mertebe mantığın geleneksel sunumlarında tanım evreninin boş
olmadığı \emph{sessizce} varsayılır. \Lean{}'de ise tipler boş olabilir
ve bu varsayım her seferinde açıkça yapılmalıdır. Bir kitaptaki ispatı
\Lean{}'e aktarırken \lc{Inhabited} ya da \lc{Nonempty} varsayımının
eksik olduğu yerler, çoğu zaman bu inceliğin gözden kaçtığı yerlerdir.
\end{dikkat}

% ---------------------------------------------------------------------
\section{Curry--Howard: bağımlı tipler}\label{sec:curry-howard-3}

Niceleyiciler, Curry--Howard tablosuna iki yeni satır ekler. Bunlar
tablonun en derin satırlarıdır, çünkü \emph{bağımlı tiplere} karşılık
gelirler (\cref{tab:curry-howard-3}).

\begin{table}[htbp]
\centering
\renewcommand{\arraystretch}{1.55}
\rowcolors{2}{tablosatir}{white}
\begin{tabular}{>{\raggedright\arraybackslash}p{2.6cm}
                >{\raggedright\arraybackslash}p{4.6cm}
                >{\raggedright\arraybackslash}p{4.4cm}}
\rowcolor{tabloust}
\tablobaslik{Mantık} & \tablobaslik{Tip kuramı} & \tablobaslik{Lean'deki kanıt} \\
$P \to Q$                & fonksiyon tipi $P \to Q$                      & \lc{fun hp => ...} \\
$\forall x{:}\alpha,\; P(x)$ & \textbf{bağımlı fonksiyon} tipi $\Pi_{x:\alpha} P(x)$ & \lc{fun x => ...} \\
$P \wedge Q$             & çarpım tipi $P \times Q$                      & \lc{⟨hp, hq⟩} \\
$\exists x{:}\alpha,\; P(x)$ & \textbf{bağımlı çift} tipi $\Sigma_{x:\alpha} P(x)$ & \lc{⟨t, ht⟩} \\
$t = s$                  & eşitlik tipi                                   & \lc{rfl} \\
\end{tabular}
\caption{Curry--Howard eşlemesinin niceleyicilere genişlemesi. $\forall$,
$\to$'nun; $\exists$ ise $\wedge$'nin bağımlı biçimidir.}
\label{tab:curry-howard-3}
\end{table}

\begin{figure}[htbp]
\centering
\begin{tikzpicture}[
    lif/.style={draw=#1, fill=#1!8, thick, rounded corners=2pt,
                minimum width=0.8cm, minimum height=1.7cm},
  ]
  % Π
  \node[font=\baslikfont\bfseries\small, text=ikincil] at (2.2,2.75) {$\Pi_{x} P(x)$: her liften bir seçim};
  \foreach \i in {0,...,4} {
    \node[lif=ikincil] (l\i) at ({1.1*\i},0.9) {};
    \node[font=\scriptsize, text=ikincil] at ({1.1*\i},2.0) {$P(\i)$};
    \node[font=\scriptsize] at ({1.1*\i},-0.2) {$\i$};
  }
  \draw[ikincil, very thick] (0,0.6) -- (1.1,1.25) -- (2.2,0.75) -- (3.3,1.3) -- (4.4,0.8);
  \foreach \i/\y in {0/0.6, 1/1.25, 2/0.75, 3/1.3, 4/0.8} \fill[ikincil] ({1.1*\i},\y) circle (2.6pt);
  \draw[black!40] (-0.5,0.0) -- (4.9,0.0);
  \node[font=\baslikfont\scriptsize, text=black!60] at (2.2,-0.65) {fonksiyon: her $x$ için bir kanıt};
  % Σ
  \node[font=\baslikfont\bfseries\small, text=vurgu] at (9.0,2.75) {$\Sigma_{x} P(x)$: bir lif ve ondan bir nokta};
  \foreach \i in {0,...,4} {
    \node[lif=vurgu, opacity=0.55] (m\i) at ({6.8+1.1*\i},0.9) {};
    \node[font=\scriptsize, text=vurgu] at ({6.8+1.1*\i},2.0) {$P(\i)$};
    \node[font=\scriptsize] at ({6.8+1.1*\i},-0.2) {$\i$};
  }
  \node[lif=vurgu, very thick] at (9.0,0.9) {};
  \fill[vurgu] (9.0,1.05) circle (3pt);
  \node[font=\baslikfont\scriptsize\bfseries, text=vurgu, anchor=west] (et) at (12.0,1.55) {$\langle 2, h\rangle$};
  \draw[vurgu, thin, -{Stealth[length=1.6mm]}] (et.west) to[bend right=15] (9.08,1.1);
  \draw[black!40] (6.3,0.0) -- (11.7,0.0);
  \node[font=\baslikfont\scriptsize, text=black!60] at (9.0,-0.65) {çift: bir tanık ve onun lifinde bir kanıt};
\end{tikzpicture}
\caption{Bağımlı tiplerin geometrik resmi. Her $x$ üzerinde bir ``lif''
$P(x)$ durur. $\Pi$ her liften bir eleman seçer; $\Sigma$ ise bir lif
ve o liften bir eleman seçer.}
\label{fig:pi-sigma}
\end{figure}

\begin{onemlifikir}[Önemli Fikir: Bağımlı tipler mantığın tamamını taşır]
Bölüm~2'de \lc{→}, \lc{∧}, \lc{∨} için başlattığımız tablo, bağımlı
tiplerle yüklemler mantığının tamamını kapsar. \Lean{}'in mantığı,
özünde \emph{bağımlı tip kuramıdır}: Yeni bir mantık kuralı
eklememize gerek kalmadan, yalnızca tipler ve fonksiyonlarla bütün
matematiği ifade edebiliriz. Bu kuramın kesin tanımını
\cref{bol:bagimli}'te vereceğiz \citep{martinlof1984, coquand1988}.
\end{onemlifikir}

% ---------------------------------------------------------------------
\section{Taktik özeti}

\begin{center}
\renewcommand{\arraystretch}{1.45}
\rowcolors{2}{tablosatir}{white}
\begin{tabularx}{\linewidth}{>{\raggedright\arraybackslash}p{3.6cm}
                             >{\raggedright\arraybackslash}p{1.5cm}
                             >{\raggedright\arraybackslash}X}
\rowcolor{tabloust}
\tablobaslik{Taktik / terim} & \tablobaslik{Kural} & \tablobaslik{Ne yapar?} \\
\leantac{intro x}                   & $\forall$G & Keyfi bir \lc{x} alır; hedef \lc{P x} olur. \\
\lc{h t}, \leantac{specialize h t}  & $\forall$E & Genel hipotezi \lc{t}'ye uygular. \\
\lc{⟨t, ht⟩}, \leantac{refine ⟨t, ?\_⟩} & $\exists$G & Tanığı \lc{t} olan bir varlık kanıtı kurar. \\
\leantac{obtain ⟨a, ha⟩ := h}       & $\exists$E & Tanığa ad verir: \lc{a} ve \lc{ha : P a}. \\
\lc{rfl}                            & $=$G       & \lc{t = t} hedefini kapatır. \\
\leantac{rw [h]}, \leantac{rw [← h]} & $=$E      & Hedefte \lc{a}'yı \lc{b}'ye (ya da tersine) çevirir. \\
\lc{h ▸ hp}                         & $=$E       & \lc{P a}'yı \lc{P b}'ye taşır. \\
\lc{h.symm}, \lc{h₁.trans h₂}       & ---        & Simetri, geçişlilik. \\
\lc{congrArg f h}                   & ---        & \lc{a = b}'den \lc{f a = f b}. \\
\leankw{calc}                       & ---        & Eşitlik ya da eşitsizlik zinciri. \\
\end{tabularx}
\end{center}

% ---------------------------------------------------------------------
\begin{arastirmanotu}
\textbf{Tanıklar Prop'tan kaçamaz.} \Lean{}'de \lc{∃ x, P x} bir
\lc{Prop}'tur ve \lc{Prop} kanıt-ilgisizdir. Bu yüzden bir varlık
kanıtından tanığı \emph{bir program olarak} çıkaramayız: \leantac{obtain}
yalnızca hedef de bir \lc{Prop} ise çalışır. Tanığı veri olarak
istediğimizde ya \lc{Type} düzeyindeki bağımlı çifti (\lc{Σ x, P x} ya da
\lc{\{x // P x\}}) kullanırız, ya da hesaplanamayan
\lc{Classical.choose}'a başvururuz. Bu tasarım, \cref{bol:yapici-klasik}'teki
``klasik akıl yürütme hesaplamayı bozmaz'' dengesinin bedelidir
\citep{carneiro2019}.

\textbf{Hilbert'in $\varepsilon$ operatörü.} Hilbert ve Ackermann'ın
mantık sisteminde $\varepsilon x.\,P(x)$ terimi, $P$'yi sağlayan
\emph{bir} nesneyi (varsa) gösterir \citep{hilbert1928}.
\lc{Classical.choose}, bu fikrin tip kuramındaki karşılığıdır.

\textbf{İleri okuma.} Birinci mertebe mantığın klasik bir sunumu için
\citet{vandalen2013}; içki içen paradoksu ve benzeri mantık bulmacaları
için \citet{smullyan1978}; tip kuramı açısından niceleyiciler için
\citet{nederpelt2014}.
\end{arastirmanotu}

\begin{bolumozeti}
\begin{itemize}
  \item Bir yüklem \lc{α → Prop} türünde bir fonksiyondur.
  \item $\forall$G: keyfi bir \lc{x} al (\leantac{intro}); $\forall$E:
        genel kanıtı bir bireye uygula. $\forall$'nın kanıtı bağımlı bir
        fonksiyondur.
  \item $\exists$G: bir tanık ve kanıtını ver (\lc{⟨t, ht⟩}); $\exists$E:
        tanığa ad ver (\leantac{obtain}). $\exists$'nin kanıtı bağımlı
        bir çifttir.
  \item $\neg\exists \leftrightarrow \forall\neg$ yapıcıdır;
        $\neg\forall \to \exists\neg$ klasiktir. $\exists\forall \to
        \forall\exists$ geçerlidir, tersi değildir.
  \item Yerine koyarken bağlı değişkenler yakalanmamalıdır.
  \item Eşitlik: \lc{rfl}, \leantac{rw}, \lc{▸}, \lc{calc}. Leibniz ilkesi
        \Lean{}'de bir teoremdir.
  \item Berber paradoksu yapıcı olarak ispatlanır. $\forall \to \exists$
        yalnızca boş olmayan evrenlerde geçerlidir.
\end{itemize}
\end{bolumozeti}

\begin{alistirmalar}
\alistirma{1} \lc{(∀ x, P x ∧ Q x) → ∀ x, P x}. (\lc{a6\_1})

\alistirma{1} \lc{(∀ x, P x) → (∀ x, Q x) → ∀ x, P x ∧ Q x}. (\lc{a6\_2})

\alistirma{1} \lc{(∃ x, P x ∧ Q x) → ∃ x, P x}. (\lc{a6\_3})

\alistirma{2} $\forall$'nın $\wedge$ üzerine dağıldığını ispatlayın:
\lc{(∀ x, P x ∧ Q x) ↔ (∀ x, P x) ∧ (∀ x, Q x)}. (\lc{a6\_4})

\alistirma{2} $\exists$'nin $\vee$ üzerine dağıldığını ispatlayın:
\lc{(∃ x, P x ∨ Q x) ↔ (∃ x, P x) ∨ (∃ x, Q x)}. (\lc{a6\_5})

\alistirma{2} \lc{(∃ x, P x ∧ Q x) → (∃ x, P x) ∧ (∃ x, Q x)}. Tersinin
neden yanlış olduğunu doğal sayılar üzerinde bir karşı örnekle
gösterin. (\lc{a6\_6})

\alistirma{2} \lc{(∀ x, P x → Q x) → (∃ x, P x) → ∃ x, Q x}. (\lc{a6\_7})

\alistirma{2} \lc{a = b} ve \lc{c = b} ise \lc{a = c}. Önce terim modunda
(\lc{symm}, \lc{trans}), sonra \leantac{rw} ile ispatlayın. (\lc{a6\_8})

\alistirma{2} \lc{a = b → f (f a) = f (f b)}. (\lc{a6\_9})

\alistirma{3} (Klasik) \lc{(¬ ∀ x, P x) ↔ ∃ x, ¬ P x}. Hangi yön klasik
mantık gerektiriyor? (\lc{a6\_10})

\alistirma{3} (Klasik) \textbf{İçki içen paradoksu}: \lc{[Inhabited α]}
varsayımıyla \lc{∃ x, P x → ∀ y, P y}. ``Bir barda öyle biri vardır ki,
o içiyorsa herkes içiyordur.'' Bu önerme neden paradoks gibi görünür?
\lc{Inhabited} varsayımı neden gereklidir? (\lc{a6\_11})

\alistirma{3} Berber paradoksunun genel hâli:
\lc{¬ ∃ x, ∀ y, R x y ↔ ¬ R y y}. Bu teoremden berber paradoksunu
tek satırda elde edin. (\lc{a6\_12})

\alistirma{2} (Kâğıt-kalem) Aşağıdaki önermelerde serbest ve bağlı
değişkenleri belirleyin; ardından $x$ yerine $y + 1$ terimini
yakalamadan koyun:
(a) $\forall y,\; x < y$; \quad (b) $(\exists x,\; P(x)) \wedge Q(x)$.

\alistirma{3} (Kâğıt-kalem) $\forall$G kuralındaki öz değişken koşulu
çiğnenirse $P(a) \to \forall x\,P(x)$ ``ispatlanabilir''. Bu yanlış
türetimi yazın ve hangi adımın kuralı çiğnediğini gösterin.
\end{alistirmalar}
